% Options for packages loaded elsewhere
\PassOptionsToPackage{unicode}{hyperref}
\PassOptionsToPackage{hyphens}{url}
\documentclass[
  11pt,
]{article}
\usepackage{xcolor}
\usepackage[margin=1in]{geometry}
\usepackage{amsmath,amssymb}
\setcounter{secnumdepth}{5}
\usepackage{iftex}
\ifPDFTeX
  \usepackage[T1]{fontenc}
  \usepackage[utf8]{inputenc}
  \usepackage{textcomp} % provide euro and other symbols
\else % if luatex or xetex
  \usepackage{unicode-math} % this also loads fontspec
  \defaultfontfeatures{Scale=MatchLowercase}
  \defaultfontfeatures[\rmfamily]{Ligatures=TeX,Scale=1}
\fi
\usepackage{lmodern}
\ifPDFTeX\else
  % xetex/luatex font selection
\fi
% Use upquote if available, for straight quotes in verbatim environments
\IfFileExists{upquote.sty}{\usepackage{upquote}}{}
\IfFileExists{microtype.sty}{% use microtype if available
  \usepackage[]{microtype}
  \UseMicrotypeSet[protrusion]{basicmath} % disable protrusion for tt fonts
}{}
\makeatletter
\@ifundefined{KOMAClassName}{% if non-KOMA class
  \IfFileExists{parskip.sty}{%
    \usepackage{parskip}
  }{% else
    \setlength{\parindent}{0pt}
    \setlength{\parskip}{6pt plus 2pt minus 1pt}}
}{% if KOMA class
  \KOMAoptions{parskip=half}}
\makeatother
\usepackage{longtable,booktabs,array}
\newcounter{none} % for unnumbered tables
\usepackage{calc} % for calculating minipage widths
% Correct order of tables after \paragraph or \subparagraph
\usepackage{etoolbox}
\makeatletter
\patchcmd\longtable{\par}{\if@noskipsec\mbox{}\fi\par}{}{}
\makeatother
% Allow footnotes in longtable head/foot
\IfFileExists{footnotehyper.sty}{\usepackage{footnotehyper}}{\usepackage{footnote}}
\makesavenoteenv{longtable}
\setlength{\emergencystretch}{3em} % prevent overfull lines
\providecommand{\tightlist}{%
  \setlength{\itemsep}{0pt}\setlength{\parskip}{0pt}}
\usepackage{bookmark}
\IfFileExists{xurl.sty}{\usepackage{xurl}}{} % add URL line breaks if available
\urlstyle{same}
\hypersetup{
  hidelinks,
  pdfcreator={LaTeX via pandoc}}

\author{}
\date{}

\begin{document}

\section{FDCAN-ABH Converter Interface Control Document
(ICD)}\label{fdcan-abh-converter-interface-control-document-icd}

\textbf{Document Number:} FDCAN-ABH-ICD-001 \textbf{Revision:} 1.0
\textbf{Date:} November 17, 2025 \textbf{Organization:} PSYONIC Inc.

\begin{center}\rule{0.5\linewidth}{0.5pt}\end{center}

\subsection{Document Control}\label{document-control}

{\def\LTcaptype{none} % do not increment counter
\begin{longtable}[]{@{}llll@{}}
\toprule\noalign{}
Version & Date & Author & Description \\
\midrule\noalign{}
\endhead
\bottomrule\noalign{}
\endlastfoot
1.0 & 2025-11-17 & Jesse Cornman & Initial release \\
\end{longtable}
}

\begin{center}\rule{0.5\linewidth}{0.5pt}\end{center}

\subsection{Table of Contents}\label{table-of-contents}

\begin{enumerate}
\def\labelenumi{\arabic{enumi}.}
\tightlist
\item
  \hyperref[1-scope]{Scope}
\item
  \hyperref[2-referenced-documents]{Referenced Documents}
\item
  \hyperref[3-interface-overview]{Interface Overview}
\item
  \hyperref[4-interface-characteristics]{Interface Characteristics}
\item
  \hyperref[5-memory-map-and-register-definitions]{Memory Map and
  Register Definitions}
\item
  \hyperref[6-operational-modes]{Operational Modes}
\item
  \hyperref[7-non-volatile-configuration-management]{Non-Volatile
  Configuration Management}
\item
  \hyperref[8-physical-and-electrical-characteristics]{Physical and
  Electrical Characteristics}
\item
  \hyperref[9-appendices]{Appendices}
\end{enumerate}

\begin{center}\rule{0.5\linewidth}{0.5pt}\end{center}

\subsection{1. Scope}\label{scope}

\subsubsection{1.1 Identification}\label{identification}

This Interface Control Document (ICD) defines the electrical,
functional, and protocol interfaces for the \textbf{FDCAN-ABH Converter}
firmware. This firmware implements a bidirectional bridge between:

\begin{itemize}
\tightlist
\item
  \textbf{CAN interface} using the DARTT (Dual Address Real-Time
  Transport) protocol
\item
  \textbf{UART interface} using the Psyonic Ability Hand Extended Mode
  API protocol
\end{itemize}

\subsubsection{1.2 System Overview}\label{system-overview}

The FDCAN-ABH Converter serves as a protocol translation gateway
enabling CAN-based controllers to command and monitor Psyonic Ability
Hand prosthetic devices. The converter operates on an STM32G4
microcontroller with FDCAN and UART peripherals.

\textbf{Primary Functions:} - Translate DARTT CAN commands to Ability
Hand UART frames - Perform HDLC byte stuffing/unstuffing for UART
communication - Provide high-level structured motor control API -
Support direct UART pass-through mode for custom protocols - Maintain
non-volatile configuration storage

\subsubsection{1.3 Document Overview}\label{document-overview}

This document specifies: - Register-level memory map using DARTT word
addressing - CAN and UART protocol requirements - Operational modes and
command sequences - Timing and performance characteristics -
Configuration procedures

\subsubsection{1.4 Intended Audience}\label{intended-audience}

This document is intended for: - Embedded systems engineers developing
CAN-based controllers - Integration engineers implementing Ability Hand
control systems - Test and validation engineers - Technical support
personnel

\begin{center}\rule{0.5\linewidth}{0.5pt}\end{center}

\subsection{2. Referenced Documents}\label{referenced-documents}

\subsubsection{2.1 Applicable Documents}\label{applicable-documents}

{\def\LTcaptype{none} % do not increment counter
\begin{longtable}[]{@{}
  >{\raggedright\arraybackslash}p{(\linewidth - 2\tabcolsep) * \real{0.5000}}
  >{\raggedright\arraybackslash}p{(\linewidth - 2\tabcolsep) * \real{0.5000}}@{}}
\toprule\noalign{}
\begin{minipage}[b]{\linewidth}\raggedright
Document
\end{minipage} & \begin{minipage}[b]{\linewidth}\raggedright
Location
\end{minipage} \\
\midrule\noalign{}
\endhead
\bottomrule\noalign{}
\endlastfoot
Ability Hand Interface Control Document &
\href{./ABILITY-HAND-ICD.pdf}{docs/ABILITY-HAND-ICD.pdf} \\
DARTT Protocol Specification &
\href{../external/dartt-protocol/docs/PROTOCOL.md}{external/dartt-protocol/docs/PROTOCOL.md} \\
\end{longtable}
}

\begin{center}\rule{0.5\linewidth}{0.5pt}\end{center}

\subsection{3. Interface Overview}\label{interface-overview}

\subsubsection{3.1 Interface Summary}\label{interface-summary}

{\def\LTcaptype{none} % do not increment counter
\begin{longtable}[]{@{}
  >{\raggedright\arraybackslash}p{(\linewidth - 6\tabcolsep) * \real{0.2683}}
  >{\raggedright\arraybackslash}p{(\linewidth - 6\tabcolsep) * \real{0.2683}}
  >{\raggedright\arraybackslash}p{(\linewidth - 6\tabcolsep) * \real{0.2439}}
  >{\raggedright\arraybackslash}p{(\linewidth - 6\tabcolsep) * \real{0.2195}}@{}}
\toprule\noalign{}
\begin{minipage}[b]{\linewidth}\raggedright
Interface
\end{minipage} & \begin{minipage}[b]{\linewidth}\raggedright
Direction
\end{minipage} & \begin{minipage}[b]{\linewidth}\raggedright
Protocol
\end{minipage} & \begin{minipage}[b]{\linewidth}\raggedright
Purpose
\end{minipage} \\
\midrule\noalign{}
\endhead
\bottomrule\noalign{}
\endlastfoot
CANH/CANL & Input/Output & DARTT & Send and receive data to control the
Ability Hand and configure the interface controller \\
UART TX & Output & Ability Hand API & Send motor commands to
prosthetic \\
UART RX & Input & Ability Hand API & Receive feedback from prosthetic \\
Flash & Storage & N/A & Persist configuration across power cycles \\
\end{longtable}
}

\subsubsection{3.2 Key Features}\label{key-features}

\begin{itemize}
\tightlist
\item
  \textbf{Dual Operating Modes:} High-level API or low-level
  pass-through
\item
  \textbf{Automatic Frame Construction:} Generate properly formatted
  UART frames from structured commands
\item
  \textbf{HDLC Processing:} Automatic unstuffing on RX, manual stuffing
  required on TX for pass-through
\item
  \textbf{Non-Volatile Storage:} Flash-based persistence for CAN and
  UART configuration
\item
  \textbf{Real-Time Performance:} Sub-millisecond latency for typical
  command-response cycles
\end{itemize}

\begin{center}\rule{0.5\linewidth}{0.5pt}\end{center}

\subsection{4. Interface
Characteristics}\label{interface-characteristics}

\subsubsection{4.1 CAN Interface}\label{can-interface}

\paragraph{4.1.1 Physical Layer}\label{physical-layer}

\begin{itemize}
\tightlist
\item
  \textbf{Bus Standard:} CAN FD (ISO 11898-1)
\item
  \textbf{Nominal Bit Rate:} 500 kbit/s (default, configurable)
\item
  \textbf{Data Bit Rate:} Same as nominal (FDCAN not utilizing data
  phase speed-up)
\item
  \textbf{Transceiver:} Integrated 5V CAN transceiver (on-board)
\item
  \textbf{Termination:} Integrated 120$\Omega$ termination resistor (on-board)
\item
  \textbf{Bus Signals:} CANH and CANL differential pair
\end{itemize}

\paragraph{4.1.2 CAN Addressing}\label{can-addressing}

The converter uses a configurable CAN arbitration ID scheme:

{\def\LTcaptype{none} % do not increment counter
\begin{longtable}[]{@{}
  >{\raggedright\arraybackslash}p{(\linewidth - 4\tabcolsep) * \real{0.2200}}
  >{\raggedright\arraybackslash}p{(\linewidth - 4\tabcolsep) * \real{0.3000}}
  >{\raggedright\arraybackslash}p{(\linewidth - 4\tabcolsep) * \real{0.4800}}@{}}
\toprule\noalign{}
\begin{minipage}[b]{\linewidth}\raggedright
Parameter
\end{minipage} & \begin{minipage}[b]{\linewidth}\raggedright
Default Value
\end{minipage} & \begin{minipage}[b]{\linewidth}\raggedright
Configuration Register
\end{minipage} \\
\midrule\noalign{}
\endhead
\bottomrule\noalign{}
\endlastfoot
Module Number & 0x50 & MODULE\_NUMBER (0x000) \\
Complementary CAN ID & 0x7FF - module\_number & Calculated (0x7AF
default) \\
\end{longtable}
}

\textbf{Addressing Rules:} - \textbf{DEFAULT CAN ID}: 0x7AF - Primary ID
- not used - Complementary ID used for all DARTT messages - Module
number must be unique on the CAN bus - Valid range: 0x01 - 0x7FE

If the MODULE\_NUMBER is re-assigned, the corresponding complementary
CAN ID will be recalculated on initialization/reset, and that ID will be
used for all DARTT messaging with the interface device.

\paragraph{4.1.3 DARTT Protocol}\label{dartt-protocol}

The DARTT (Dual Address Real-Time Transport) protocol provides block
memory read/write access over CAN.

\textbf{Key Characteristics:} - \textbf{Word Size:} 32 bits (4 bytes) -
\textbf{Addressing:} Word-aligned (addresses are word indexes, not byte
offsets) - \textbf{Access Modes:} Block read, block write -
\textbf{Endianness:} Little-endian (index and num\_bytes fields) -
\textbf{Payload Endianness:} Application-defined - Little-endian for
this device - \textbf{Maximum Block Size:} Application-defined - 6 bytes
maximum for this device - \textbf{Message Type:}
TYPE\_ADDR\_CRC\_MESSAGE (Type 2) - CAN provides addressing and error
checking

\textbf{CAN Frame Formats:}

This device uses DARTT Type 2 messages (TYPE\_ADDR\_CRC\_MESSAGE), which
rely on CAN's built-in arbitration ID and CRC, eliminating protocol
overhead.

\textbf{Write Frame (Controller → Converter):}

{\def\LTcaptype{none} % do not increment counter
\begin{longtable}[]{@{}ll@{}}
\toprule\noalign{}
Bytes 0-1 & Bytes 2-N \\
\midrule\noalign{}
\endhead
\bottomrule\noalign{}
\endlastfoot
Index (R=0) & Payload Data \\
\end{longtable}
}

\begin{itemize}
\tightlist
\item
  \textbf{Index}: 16-bit little-endian word index (bit 15 = 0 for write)
\item
  \textbf{Payload}: N bytes of data to write (between 4 and 6 on this
  device for standard CAN compatibility)
\end{itemize}

\textbf{Read Request Frame (Controller → Converter):}

{\def\LTcaptype{none} % do not increment counter
\begin{longtable}[]{@{}ll@{}}
\toprule\noalign{}
Bytes 0-1 & Bytes 2-3 \\
\midrule\noalign{}
\endhead
\bottomrule\noalign{}
\endlastfoot
Index (R=1) & Num Bytes \\
\end{longtable}
}

\begin{itemize}
\tightlist
\item
  \textbf{Index}: 16-bit little-endian word index (bit 15 = 1 for read)
\item
  \textbf{Num Bytes}: 16-bit little-endian byte count to read.
  Requesting more than 8 bytes from this device will result in no
  response.
\end{itemize}

\textbf{Read Reply Frame (Converter → Controller):}

{\def\LTcaptype{none} % do not increment counter
\begin{longtable}[]{@{}l@{}}
\toprule\noalign{}
Bytes 0-N \\
\midrule\noalign{}
\endhead
\bottomrule\noalign{}
\endlastfoot
Requested Data Block \\
\end{longtable}
}

\begin{itemize}
\tightlist
\item
  \textbf{Data Block}: N bytes of requested data (raw payload)
\end{itemize}

\subsubsection{4.2 UART Interface}\label{uart-interface}

\paragraph{4.2.1 Physical Layer}\label{physical-layer-1}

\begin{itemize}
\tightlist
\item
  \textbf{Baud Rate:} 460800 bps (default, configurable)
\item
  \textbf{Data Format:} 8 data bits, no parity, 1 stop bit (8N1)
\item
  \textbf{Flow Control:} None
\item
  \textbf{Direction:} Full duplex
\end{itemize}

\paragraph{4.2.2 Supported Baud Rates}\label{supported-baud-rates}

{\def\LTcaptype{none} % do not increment counter
\begin{longtable}[]{@{}lll@{}}
\toprule\noalign{}
Baud Rate (bps) & UART\_BAUD\_RATE Value & Common Use \\
\midrule\noalign{}
\endhead
\bottomrule\noalign{}
\endlastfoot
9600 & 9600 & Debugging, legacy \\
19200 & 19200 & Low-speed reliable \\
38400 & 38400 & Standard \\
57600 & 57600 & Standard \\
115200 & 115200 & Common high-speed \\
230400 & 230400 & High-speed \\
460800 & 460800 & Default (recommended) \\
921600 & 921600 & Maximum reliable \\
1000000 & 1000000 & Maximum \\
\end{longtable}
}

\textbf{Note:} Baud rates other than 460800 require modification to the
baud rate setting on the Ability Hand. Additionally, when using baud
rates over 460800, it is recommended to disable BLE, as bluetooth
activity at high baudrates may result in dropped frames.

\paragraph{4.2.3 HDLC Framing}\label{hdlc-framing}

The UART frame synchronization is based on HDLC (High-Level Data Link
Control) byte stuffing. It deviates from this specification slightly in
that a unique frame delimiter is expected at both the beginning and end
of each frame. Refer to the
\href{https://github.com/ocanath/byte-stuffing}{following library} for
compatible C and Python implementations.

\textbf{Stuffing Requirements (TX Pass-Through Mode):} - Controller MUST
perform stuffing before writing to UART\_TX\_MEM - Failure to stuff will
result in no communication to/from the hand

\paragraph{4.2.4 Ability Hand Protocol
Summary}\label{ability-hand-protocol-summary}

The downstream UART protocol supports multiple control modes:

{\def\LTcaptype{none} % do not increment counter
\begin{longtable}[]{@{}
  >{\raggedright\arraybackslash}p{(\linewidth - 6\tabcolsep) * \real{0.2857}}
  >{\raggedright\arraybackslash}p{(\linewidth - 6\tabcolsep) * \real{0.2500}}
  >{\raggedright\arraybackslash}p{(\linewidth - 6\tabcolsep) * \real{0.1964}}
  >{\raggedright\arraybackslash}p{(\linewidth - 6\tabcolsep) * \real{0.2679}}@{}}
\toprule\noalign{}
\begin{minipage}[b]{\linewidth}\raggedright
Command Header
\end{minipage} & \begin{minipage}[b]{\linewidth}\raggedright
Control Mode
\end{minipage} & \begin{minipage}[b]{\linewidth}\raggedright
Data Type
\end{minipage} & \begin{minipage}[b]{\linewidth}\raggedright
Reply Variant
\end{minipage} \\
\midrule\noalign{}
\endhead
\bottomrule\noalign{}
\endlastfoot
0x10-0x12 & Position & int16\_t (±32767 = 0-150°) & 1, 2, or 3 \\
0x20-0x22 & Velocity & int16\_t (±32767 = 0-3000°/s) & 1, 2, or 3 \\
0x30-0x32 & Torque/Current & int16\_t & 1, 2, or 3 \\
0x40-0x42 & Voltage & int16\_t & 1, 2, or 3 \\
0xA0-0xA2 & Read-only & N/A & 1, 2, or 3 \\
\end{longtable}
}

\textbf{Reply Variants:} - \textbf{Variant 1 (TX1):} Position + Current
feedback - \textbf{Variant 2 (TX2):} Position + Velocity feedback -
\textbf{Variant 3 (TX3):} Position + Current + FSR data

\textbf{Frame Structure:}

\begin{verbatim}
[0x7E] [Address] [Command Header] [12-byte Payload] [Checksum] [0x7E]
\end{verbatim}

\textbf{Checksum:} 8-bit two's complement negation of sum of all bytes
(address through payload).

Refer to ABILITY-HAND-ICD.pdf for complete Ability Hand protocol
details.

\begin{center}\rule{0.5\linewidth}{0.5pt}\end{center}

\subsection{5. Memory Map and Register
Definitions}\label{memory-map-and-register-definitions}

\subsubsection{5.1 Memory Map Overview}\label{memory-map-overview}

The FDCAN-ABH Converter exposes all configurable parameters and
operational data through a DARTT-accessible memory map. All registers
are organized as 32-bit words with the following access types:

\begin{itemize}
\tightlist
\item
  \textbf{R/W}: Read/Write - Can be read and modified
\item
  \textbf{RO}: Read-Only - Can only be read
\item
  \textbf{WO}: Write-Only - Writing triggers an action
\item
  \textbf{R/W-NV}: Read/Write Non-Volatile - Persists across power
  cycles when committed
\end{itemize}

\paragraph{5.1.1 Memory Map Structure}\label{memory-map-structure}

{\def\LTcaptype{none} % do not increment counter
\begin{longtable}[]{@{}
  >{\raggedright\arraybackslash}p{(\linewidth - 6\tabcolsep) * \real{0.3036}}
  >{\raggedright\arraybackslash}p{(\linewidth - 6\tabcolsep) * \real{0.2143}}
  >{\raggedright\arraybackslash}p{(\linewidth - 6\tabcolsep) * \real{0.2500}}
  >{\raggedright\arraybackslash}p{(\linewidth - 6\tabcolsep) * \real{0.2321}}@{}}
\toprule\noalign{}
\begin{minipage}[b]{\linewidth}\raggedright
Base Word Index
\end{minipage} & \begin{minipage}[b]{\linewidth}\raggedright
Block Name
\end{minipage} & \begin{minipage}[b]{\linewidth}\raggedright
Size (words)
\end{minipage} & \begin{minipage}[b]{\linewidth}\raggedright
Description
\end{minipage} \\
\midrule\noalign{}
\endhead
\bottomrule\noalign{}
\endlastfoot
0x000 & Non-Volatile Configuration & 6 & CAN and UART configuration
parameters \\
0x006 & Ability Hand Control API & 40 & High-level motor control
interface \\
0x02E & Register Access Interface & 3 & Direct register read/write
(unimplemented) \\
0x031 & ABH Read Timeout & 1 & Reply timeout configuration \\
0x032 & UART RX Decoded Buffer & 19 & HDLC-decoded receive buffer (76
bytes) \\
0x045 & UART RX Byte Count & 1 & Valid bytes in RX buffer \\
0x046 & UART TX Buffer & 19 & Raw transmit buffer. Cleared on interface
device once tranmission begins \\
0x059 & UART TX Byte Count & 1 & TX trigger - write byte count to
initiate transmission \\
0x05A & Git Hash & 4 & Zero-terminated short git hash (max 15
characters, null-terminated) - for version tracking \\
0x5E & Update nonvolatile storage & 1 & Flag for updating non-volatile
storage - cleared upon completion of update operation \\
\end{longtable}
}

\textbf{Total Memory Map Size:} 95 words (380 bytes)

\begin{center}\rule{0.5\linewidth}{0.5pt}\end{center}

\subsubsection{5.2 Block 1: Non-Volatile Configuration
(0x000-0x005)}\label{block-1-non-volatile-configuration-0x000-0x005}

These registers define persistent configuration parameters stored in
flash memory (Page 63). Changes to these registers require setting the
\texttt{UPDATE\_NONVOLATILE\_STORAGE} flag (0x057) to persist across
power cycles.

{\def\LTcaptype{none} % do not increment counter
\begin{longtable}[]{@{}
  >{\raggedright\arraybackslash}p{(\linewidth - 10\tabcolsep) * \real{0.1905}}
  >{\raggedright\arraybackslash}p{(\linewidth - 10\tabcolsep) * \real{0.2381}}
  >{\raggedright\arraybackslash}p{(\linewidth - 10\tabcolsep) * \real{0.0952}}
  >{\raggedright\arraybackslash}p{(\linewidth - 10\tabcolsep) * \real{0.1270}}
  >{\raggedright\arraybackslash}p{(\linewidth - 10\tabcolsep) * \real{0.1429}}
  >{\raggedright\arraybackslash}p{(\linewidth - 10\tabcolsep) * \real{0.2063}}@{}}
\toprule\noalign{}
\begin{minipage}[b]{\linewidth}\raggedright
Word Index
\end{minipage} & \begin{minipage}[b]{\linewidth}\raggedright
Register Name
\end{minipage} & \begin{minipage}[b]{\linewidth}\raggedright
Type
\end{minipage} & \begin{minipage}[b]{\linewidth}\raggedright
Access
\end{minipage} & \begin{minipage}[b]{\linewidth}\raggedright
Default
\end{minipage} & \begin{minipage}[b]{\linewidth}\raggedright
Description
\end{minipage} \\
\midrule\noalign{}
\endhead
\bottomrule\noalign{}
\endlastfoot
0x000 & MODULE\_NUMBER & uint32\_t & R/W-NV & 0x50 & CAN module address.
Complementary address: 0x7FF - module\_number \\
0x001 & UART\_BAUD\_RATE & uint32\_t & R/W-NV & 460800 & UART baud rate
in bits/second. See Section 4.2.2 for valid values \\
0x002 & FDCAN\_NBRP & uint32\_t & R/W-NV & 2 & FDCAN Nominal Bit Rate
Prescaler. Range: 1-512 \\
0x003 & FDCAN\_NTSEG1 & uint32\_t & R/W-NV & 135 & FDCAN Nominal Time
Segment 1. Range: 2-256 \\
0x004 & FDCAN\_NTSEG2 & uint32\_t & R/W-NV & 34 & FDCAN Nominal Time
Segment 2. Range: 2-128 \\
0x005 & UNUSED\_ZEROPAD & uint32\_t & RO & 0 & Unused nonvolatile
storage word. \\
\end{longtable}
}

\textbf{CAN Bit Rate Calculation:}

\begin{verbatim}
Nominal_Bit_Rate = 170_MHz / (NBRP &times; (1 + NTSEG1 + NTSEG2))
\end{verbatim}

\textbf{Default Example:}

\begin{verbatim}
170_MHz / (2 &times; (1 + 135 + 34)) = 170_MHz / 340 = 500 kbit/s
\end{verbatim}

\textbf{Configuration Notes:} - Changes take effect after power cycle or
system reset - Invalid CAN timing parameters may cause bus communication
failure - Peripheral clock (PCLK): 170 MHz (STM32G4 maximum)

\begin{center}\rule{0.5\linewidth}{0.5pt}\end{center}

\subsubsection{5.3 Block 2: Ability Hand Control API
(0x006-0x026)}\label{block-2-ability-hand-control-api-0x006-0x026}

This block provides high-level structured control of the Ability Hand.
Writing to the command setpoint arrays automatically generates properly
formatted UART frames with HDLC stuffing. Reply data is parsed and
populated into the feedback arrays.

\paragraph{5.3.1 Configuration and Control
(0x006-0x007)}\label{configuration-and-control-0x006-0x007}

{\def\LTcaptype{none} % do not increment counter
\begin{longtable}[]{@{}
  >{\raggedright\arraybackslash}p{(\linewidth - 10\tabcolsep) * \real{0.1905}}
  >{\raggedright\arraybackslash}p{(\linewidth - 10\tabcolsep) * \real{0.2381}}
  >{\raggedright\arraybackslash}p{(\linewidth - 10\tabcolsep) * \real{0.0952}}
  >{\raggedright\arraybackslash}p{(\linewidth - 10\tabcolsep) * \real{0.1270}}
  >{\raggedright\arraybackslash}p{(\linewidth - 10\tabcolsep) * \real{0.1429}}
  >{\raggedright\arraybackslash}p{(\linewidth - 10\tabcolsep) * \real{0.2063}}@{}}
\toprule\noalign{}
\begin{minipage}[b]{\linewidth}\raggedright
Word Index
\end{minipage} & \begin{minipage}[b]{\linewidth}\raggedright
Register Name
\end{minipage} & \begin{minipage}[b]{\linewidth}\raggedright
Type
\end{minipage} & \begin{minipage}[b]{\linewidth}\raggedright
Access
\end{minipage} & \begin{minipage}[b]{\linewidth}\raggedright
Default
\end{minipage} & \begin{minipage}[b]{\linewidth}\raggedright
Description
\end{minipage} \\
\midrule\noalign{}
\endhead
\bottomrule\noalign{}
\endlastfoot
0x006 & ABH\_ADDRESS & uint32\_t & R/W & 0x50 & Ability Hand UART device
address \\
0x007 & ABH\_COMMAND\_HEADER & uint32\_t & R/W & 0x10 & Frame command
header. See Section 5.3.2 \\
\end{longtable}
}

\paragraph{5.3.2 Command Header Values}\label{command-header-values}

{\def\LTcaptype{none} % do not increment counter
\begin{longtable}[]{@{}
  >{\raggedright\arraybackslash}p{(\linewidth - 6\tabcolsep) * \real{0.1750}}
  >{\raggedright\arraybackslash}p{(\linewidth - 6\tabcolsep) * \real{0.1500}}
  >{\raggedright\arraybackslash}p{(\linewidth - 6\tabcolsep) * \real{0.3250}}
  >{\raggedright\arraybackslash}p{(\linewidth - 6\tabcolsep) * \real{0.3500}}@{}}
\toprule\noalign{}
\begin{minipage}[b]{\linewidth}\raggedright
Value
\end{minipage} & \begin{minipage}[b]{\linewidth}\raggedright
Mode
\end{minipage} & \begin{minipage}[b]{\linewidth}\raggedright
Description
\end{minipage} & \begin{minipage}[b]{\linewidth}\raggedright
Reply Format
\end{minipage} \\
\midrule\noalign{}
\endhead
\bottomrule\noalign{}
\endlastfoot
0xA0 & Read-only & Dummy command, reply variant 1 & Position +
Current \\
0xA1 & Read-only & Dummy command, reply variant 2 & Position +
Velocity \\
0xA2 & Read-only & Dummy command, reply variant 3 & Position + Current +
FSR \\
0x10 & Position control & Position control, reply variant 1 & Position +
Current \\
0x11 & Position control & Position control, reply variant 2 & Position +
Velocity \\
0x12 & Position control & Position control, reply variant 3 & Position +
Current + FSR \\
0x20 & Velocity control & Velocity control, reply variant 1 & Position +
Current \\
0x21 & Velocity control & Velocity control, reply variant 2 & Position +
Velocity \\
0x22 & Velocity control & Velocity control, reply variant 3 & Position +
Current + FSR \\
0x30 & Torque control & Torque control, reply variant 1 & Position +
Current \\
0x31 & Torque control & Torque control, reply variant 2 & Position +
Velocity \\
0x32 & Torque control & Torque control, reply variant 3 & Position +
Current + FSR \\
0x40 & Voltage control & Voltage control, reply variant 1 & Position +
Current \\
0x41 & Voltage control & Voltage control, reply variant 2 & Position +
Velocity \\
0x42 & Voltage control & Voltage control, reply variant 3 & Position +
Current + FSR \\
0xDE & Register write & Write to internal register & None \\
0xDA & Register read & Read from internal register & 32-bit register
value \\
0x07 & Bluetooth & Enable Bluetooth radio & None \\
0x08 & Bluetooth & Disable Bluetooth radio & None \\
0x09 & System & Restart hand controller & None \\
0x7C & API & Exit API mode & None \\
\end{longtable}
}

\paragraph{5.3.3 Motor Setpoints - Position
(0x008-0x00A)}\label{motor-setpoints---position-0x008-0x00a}

{\def\LTcaptype{none} % do not increment counter
\begin{longtable}[]{@{}
  >{\raggedright\arraybackslash}p{(\linewidth - 10\tabcolsep) * \real{0.1905}}
  >{\raggedright\arraybackslash}p{(\linewidth - 10\tabcolsep) * \real{0.2381}}
  >{\raggedright\arraybackslash}p{(\linewidth - 10\tabcolsep) * \real{0.0952}}
  >{\raggedright\arraybackslash}p{(\linewidth - 10\tabcolsep) * \real{0.1270}}
  >{\raggedright\arraybackslash}p{(\linewidth - 10\tabcolsep) * \real{0.1429}}
  >{\raggedright\arraybackslash}p{(\linewidth - 10\tabcolsep) * \real{0.2063}}@{}}
\toprule\noalign{}
\begin{minipage}[b]{\linewidth}\raggedright
Word Index
\end{minipage} & \begin{minipage}[b]{\linewidth}\raggedright
Register Name
\end{minipage} & \begin{minipage}[b]{\linewidth}\raggedright
Type
\end{minipage} & \begin{minipage}[b]{\linewidth}\raggedright
Access
\end{minipage} & \begin{minipage}[b]{\linewidth}\raggedright
Default
\end{minipage} & \begin{minipage}[b]{\linewidth}\raggedright
Description
\end{minipage} \\
\midrule\noalign{}
\endhead
\bottomrule\noalign{}
\endlastfoot
0x008 & Q\_DESIRED\_0\_1 & int16\_t{[}2{]} & R/W & 0 & Position
setpoints for motors 0-1 \\
0x009 & Q\_DESIRED\_2\_3 & int16\_t{[}2{]} & R/W & 0 & Position
setpoints for motors 2-3 \\
0x00A & Q\_DESIRED\_4\_5 & int16\_t{[}2{]} & R/W & 0 & Position
setpoints for motors 4-5 \\
\end{longtable}
}

\textbf{Data Format:} Each 32-bit word contains two packed int16\_t
values: - Bits {[}15:0{]}: Motor N - Bits {[}31:16{]}: Motor N+1

\textbf{Scaling:}

\begin{verbatim}
position_digital = (angle_degrees / 150.0) &times; 32767
\end{verbatim}

\paragraph{5.3.4 Motor Setpoints - Voltage
(0x00B-0x00D)}\label{motor-setpoints---voltage-0x00b-0x00d}

{\def\LTcaptype{none} % do not increment counter
\begin{longtable}[]{@{}
  >{\raggedright\arraybackslash}p{(\linewidth - 10\tabcolsep) * \real{0.1905}}
  >{\raggedright\arraybackslash}p{(\linewidth - 10\tabcolsep) * \real{0.2381}}
  >{\raggedright\arraybackslash}p{(\linewidth - 10\tabcolsep) * \real{0.0952}}
  >{\raggedright\arraybackslash}p{(\linewidth - 10\tabcolsep) * \real{0.1270}}
  >{\raggedright\arraybackslash}p{(\linewidth - 10\tabcolsep) * \real{0.1429}}
  >{\raggedright\arraybackslash}p{(\linewidth - 10\tabcolsep) * \real{0.2063}}@{}}
\toprule\noalign{}
\begin{minipage}[b]{\linewidth}\raggedright
Word Index
\end{minipage} & \begin{minipage}[b]{\linewidth}\raggedright
Register Name
\end{minipage} & \begin{minipage}[b]{\linewidth}\raggedright
Type
\end{minipage} & \begin{minipage}[b]{\linewidth}\raggedright
Access
\end{minipage} & \begin{minipage}[b]{\linewidth}\raggedright
Default
\end{minipage} & \begin{minipage}[b]{\linewidth}\raggedright
Description
\end{minipage} \\
\midrule\noalign{}
\endhead
\bottomrule\noalign{}
\endlastfoot
0x00B & VQ\_DESIRED\_0\_1 & int16\_t{[}2{]} & R/W & 0 & Voltage
setpoints for motors 0-1 \\
0x00C & VQ\_DESIRED\_2\_3 & int16\_t{[}2{]} & R/W & 0 & Voltage
setpoints for motors 2-3 \\
0x00D & VQ\_DESIRED\_4\_5 & int16\_t{[}2{]} & R/W & 0 & Voltage
setpoints for motors 4-5 \\
\end{longtable}
}

\textbf{Data Format:} Each word contains two packed int16\_t values
(motor pairs).

\textbf{Scaling:} ±32767 represents full-scale PWM duty cycle (direct
voltage control).

\textbf{Use Case:} Open-loop voltage control for testing or specialized
control algorithms.

\paragraph{5.3.5 Motor Setpoints - Current/Torque
(0x00E-0x010)}\label{motor-setpoints---currenttorque-0x00e-0x010}

{\def\LTcaptype{none} % do not increment counter
\begin{longtable}[]{@{}
  >{\raggedright\arraybackslash}p{(\linewidth - 10\tabcolsep) * \real{0.1905}}
  >{\raggedright\arraybackslash}p{(\linewidth - 10\tabcolsep) * \real{0.2381}}
  >{\raggedright\arraybackslash}p{(\linewidth - 10\tabcolsep) * \real{0.0952}}
  >{\raggedright\arraybackslash}p{(\linewidth - 10\tabcolsep) * \real{0.1270}}
  >{\raggedright\arraybackslash}p{(\linewidth - 10\tabcolsep) * \real{0.1429}}
  >{\raggedright\arraybackslash}p{(\linewidth - 10\tabcolsep) * \real{0.2063}}@{}}
\toprule\noalign{}
\begin{minipage}[b]{\linewidth}\raggedright
Word Index
\end{minipage} & \begin{minipage}[b]{\linewidth}\raggedright
Register Name
\end{minipage} & \begin{minipage}[b]{\linewidth}\raggedright
Type
\end{minipage} & \begin{minipage}[b]{\linewidth}\raggedright
Access
\end{minipage} & \begin{minipage}[b]{\linewidth}\raggedright
Default
\end{minipage} & \begin{minipage}[b]{\linewidth}\raggedright
Description
\end{minipage} \\
\midrule\noalign{}
\endhead
\bottomrule\noalign{}
\endlastfoot
0x00E & IQ\_DESIRED\_0\_1 & int16\_t{[}2{]} & R/W & 0 & Current/torque
setpoints for motors 0-1 \\
0x00F & IQ\_DESIRED\_2\_3 & int16\_t{[}2{]} & R/W & 0 & Current/torque
setpoints for motors 2-3 \\
0x010 & IQ\_DESIRED\_4\_5 & int16\_t{[}2{]} & R/W & 0 & Current/torque
setpoints for motors 4-5 \\
\end{longtable}
}

\textbf{Data Format:} Each word contains two packed int16\_t values.

\textbf{Scaling:} Hand-specific, refer to Ability Hand ICD Section 3.4
for current scaling factors.

\textbf{Use Case:} Torque control mode for force-limited grasping.

\paragraph{5.3.6 Motor Setpoints - Velocity
(0x011-0x013)}\label{motor-setpoints---velocity-0x011-0x013}

{\def\LTcaptype{none} % do not increment counter
\begin{longtable}[]{@{}
  >{\raggedright\arraybackslash}p{(\linewidth - 10\tabcolsep) * \real{0.1905}}
  >{\raggedright\arraybackslash}p{(\linewidth - 10\tabcolsep) * \real{0.2381}}
  >{\raggedright\arraybackslash}p{(\linewidth - 10\tabcolsep) * \real{0.0952}}
  >{\raggedright\arraybackslash}p{(\linewidth - 10\tabcolsep) * \real{0.1270}}
  >{\raggedright\arraybackslash}p{(\linewidth - 10\tabcolsep) * \real{0.1429}}
  >{\raggedright\arraybackslash}p{(\linewidth - 10\tabcolsep) * \real{0.2063}}@{}}
\toprule\noalign{}
\begin{minipage}[b]{\linewidth}\raggedright
Word Index
\end{minipage} & \begin{minipage}[b]{\linewidth}\raggedright
Register Name
\end{minipage} & \begin{minipage}[b]{\linewidth}\raggedright
Type
\end{minipage} & \begin{minipage}[b]{\linewidth}\raggedright
Access
\end{minipage} & \begin{minipage}[b]{\linewidth}\raggedright
Default
\end{minipage} & \begin{minipage}[b]{\linewidth}\raggedright
Description
\end{minipage} \\
\midrule\noalign{}
\endhead
\bottomrule\noalign{}
\endlastfoot
0x011 & VELOCITY\_DESIRED\_0\_1 & int16\_t{[}2{]} & R/W & 0 & Velocity
setpoints for motors 0-1 \\
0x012 & VELOCITY\_DESIRED\_2\_3 & int16\_t{[}2{]} & R/W & 0 & Velocity
setpoints for motors 2-3 \\
0x013 & VELOCITY\_DESIRED\_4\_5 & int16\_t{[}2{]} & R/W & 0 & Velocity
setpoints for motors 4-5 \\
\end{longtable}
}

\textbf{Data Format:} Each word contains two packed int16\_t values.

\textbf{Scaling (to hand):}

\begin{verbatim}
velocity_digital = (velocity_deg_per_sec / 3000.0) &times; 32767
\end{verbatim}

\textbf{Valid Range:} - 0 °/s = 0 - 1500 °/s = 16384 - 3000 °/s = 32767

\paragraph{5.3.7 Motor Feedback - Position
(0x014-0x016)}\label{motor-feedback---position-0x014-0x016}

{\def\LTcaptype{none} % do not increment counter
\begin{longtable}[]{@{}
  >{\raggedright\arraybackslash}p{(\linewidth - 10\tabcolsep) * \real{0.1905}}
  >{\raggedright\arraybackslash}p{(\linewidth - 10\tabcolsep) * \real{0.2381}}
  >{\raggedright\arraybackslash}p{(\linewidth - 10\tabcolsep) * \real{0.0952}}
  >{\raggedright\arraybackslash}p{(\linewidth - 10\tabcolsep) * \real{0.1270}}
  >{\raggedright\arraybackslash}p{(\linewidth - 10\tabcolsep) * \real{0.1429}}
  >{\raggedright\arraybackslash}p{(\linewidth - 10\tabcolsep) * \real{0.2063}}@{}}
\toprule\noalign{}
\begin{minipage}[b]{\linewidth}\raggedright
Word Index
\end{minipage} & \begin{minipage}[b]{\linewidth}\raggedright
Register Name
\end{minipage} & \begin{minipage}[b]{\linewidth}\raggedright
Type
\end{minipage} & \begin{minipage}[b]{\linewidth}\raggedright
Access
\end{minipage} & \begin{minipage}[b]{\linewidth}\raggedright
Default
\end{minipage} & \begin{minipage}[b]{\linewidth}\raggedright
Description
\end{minipage} \\
\midrule\noalign{}
\endhead
\bottomrule\noalign{}
\endlastfoot
0x014 & Q\_ACTUAL\_0\_1 & int16\_t{[}2{]} & RO & 0 & Actual position
feedback for motors 0-1 \\
0x015 & Q\_ACTUAL\_2\_3 & int16\_t{[}2{]} & RO & 0 & Actual position
feedback for motors 2-3 \\
0x016 & Q\_ACTUAL\_4\_5 & int16\_t{[}2{]} & RO & 0 & Actual position
feedback for motors 4-5 \\
\end{longtable}
}

\textbf{Scaling:} Same as position setpoints (±32767 = 0-150°).

\textbf{Update:} Populated after receiving valid UART reply from hand.
Available in all reply variants.

\paragraph{5.3.8 Motor Feedback - Current
(0x017-0x019)}\label{motor-feedback---current-0x017-0x019}

{\def\LTcaptype{none} % do not increment counter
\begin{longtable}[]{@{}
  >{\raggedright\arraybackslash}p{(\linewidth - 10\tabcolsep) * \real{0.1905}}
  >{\raggedright\arraybackslash}p{(\linewidth - 10\tabcolsep) * \real{0.2381}}
  >{\raggedright\arraybackslash}p{(\linewidth - 10\tabcolsep) * \real{0.0952}}
  >{\raggedright\arraybackslash}p{(\linewidth - 10\tabcolsep) * \real{0.1270}}
  >{\raggedright\arraybackslash}p{(\linewidth - 10\tabcolsep) * \real{0.1429}}
  >{\raggedright\arraybackslash}p{(\linewidth - 10\tabcolsep) * \real{0.2063}}@{}}
\toprule\noalign{}
\begin{minipage}[b]{\linewidth}\raggedright
Word Index
\end{minipage} & \begin{minipage}[b]{\linewidth}\raggedright
Register Name
\end{minipage} & \begin{minipage}[b]{\linewidth}\raggedright
Type
\end{minipage} & \begin{minipage}[b]{\linewidth}\raggedright
Access
\end{minipage} & \begin{minipage}[b]{\linewidth}\raggedright
Default
\end{minipage} & \begin{minipage}[b]{\linewidth}\raggedright
Description
\end{minipage} \\
\midrule\noalign{}
\endhead
\bottomrule\noalign{}
\endlastfoot
0x017 & IQ\_ACTUAL\_0\_1 & int16\_t{[}2{]} & RO & 0 & Actual current
feedback for motors 0-1 \\
0x018 & IQ\_ACTUAL\_2\_3 & int16\_t{[}2{]} & RO & 0 & Actual current
feedback for motors 2-3 \\
0x019 & IQ\_ACTUAL\_4\_5 & int16\_t{[}2{]} & RO & 0 & Actual current
feedback for motors 4-5 \\
\end{longtable}
}

\textbf{Availability:} Populated when reply variant 1 (TX1) or 3 (TX3)
is selected via command header.

\textbf{Use Case:} Current monitoring for torque estimation and fault
detection.

\paragraph{5.3.9 Motor Feedback - Velocity
(0x01A-0x01C)}\label{motor-feedback---velocity-0x01a-0x01c}

{\def\LTcaptype{none} % do not increment counter
\begin{longtable}[]{@{}
  >{\raggedright\arraybackslash}p{(\linewidth - 10\tabcolsep) * \real{0.1905}}
  >{\raggedright\arraybackslash}p{(\linewidth - 10\tabcolsep) * \real{0.2381}}
  >{\raggedright\arraybackslash}p{(\linewidth - 10\tabcolsep) * \real{0.0952}}
  >{\raggedright\arraybackslash}p{(\linewidth - 10\tabcolsep) * \real{0.1270}}
  >{\raggedright\arraybackslash}p{(\linewidth - 10\tabcolsep) * \real{0.1429}}
  >{\raggedright\arraybackslash}p{(\linewidth - 10\tabcolsep) * \real{0.2063}}@{}}
\toprule\noalign{}
\begin{minipage}[b]{\linewidth}\raggedright
Word Index
\end{minipage} & \begin{minipage}[b]{\linewidth}\raggedright
Register Name
\end{minipage} & \begin{minipage}[b]{\linewidth}\raggedright
Type
\end{minipage} & \begin{minipage}[b]{\linewidth}\raggedright
Access
\end{minipage} & \begin{minipage}[b]{\linewidth}\raggedright
Default
\end{minipage} & \begin{minipage}[b]{\linewidth}\raggedright
Description
\end{minipage} \\
\midrule\noalign{}
\endhead
\bottomrule\noalign{}
\endlastfoot
0x01A & VELOCITY\_ACTUAL\_0\_1 & int16\_t{[}2{]} & RO & 0 & Actual
velocity feedback for motors 0-1 \\
0x01B & VELOCITY\_ACTUAL\_2\_3 & int16\_t{[}2{]} & RO & 0 & Actual
velocity feedback for motors 2-3 \\
0x01C & VELOCITY\_ACTUAL\_4\_5 & int16\_t{[}2{]} & RO & 0 & Actual
velocity feedback for motors 4-5 \\
\end{longtable}
}

\textbf{Availability:} Populated when reply variant 2 (TX2) is selected
via command header.

\textbf{Scaling (from hand):}

\begin{verbatim}
velocity_rad_per_sec = velocity_digital / 4
\end{verbatim}

\textbf{Note:} Velocity feedback uses radians/sec, while velocity
setpoints use degrees/sec.

\paragraph{5.3.10 Force Sensor Feedback
(0x01D-0x024)}\label{force-sensor-feedback-0x01d-0x024}

{\def\LTcaptype{none} % do not increment counter
\begin{longtable}[]{@{}
  >{\raggedright\arraybackslash}p{(\linewidth - 10\tabcolsep) * \real{0.1905}}
  >{\raggedright\arraybackslash}p{(\linewidth - 10\tabcolsep) * \real{0.2381}}
  >{\raggedright\arraybackslash}p{(\linewidth - 10\tabcolsep) * \real{0.0952}}
  >{\raggedright\arraybackslash}p{(\linewidth - 10\tabcolsep) * \real{0.1270}}
  >{\raggedright\arraybackslash}p{(\linewidth - 10\tabcolsep) * \real{0.1429}}
  >{\raggedright\arraybackslash}p{(\linewidth - 10\tabcolsep) * \real{0.2063}}@{}}
\toprule\noalign{}
\begin{minipage}[b]{\linewidth}\raggedright
Word Index
\end{minipage} & \begin{minipage}[b]{\linewidth}\raggedright
Register Name
\end{minipage} & \begin{minipage}[b]{\linewidth}\raggedright
Type
\end{minipage} & \begin{minipage}[b]{\linewidth}\raggedright
Access
\end{minipage} & \begin{minipage}[b]{\linewidth}\raggedright
Default
\end{minipage} & \begin{minipage}[b]{\linewidth}\raggedright
Description
\end{minipage} \\
\midrule\noalign{}
\endhead
\bottomrule\noalign{}
\endlastfoot
0x01D & FSR\_RAW\_0\_1 & uint16\_t{[}2{]} & RO & 0 & FSR sensors 0-1
(unpacked values) \\
0x01E & FSR\_RAW\_2\_3 & uint16\_t{[}2{]} & RO & 0 & FSR sensors 2-3 \\
0x01F & FSR\_RAW\_4\_5 & uint16\_t{[}2{]} & RO & 0 & FSR sensors 4-5 \\
0x020 & FSR\_RAW\_6\_7 & uint16\_t{[}2{]} & RO & 0 & FSR sensors 6-7 \\
0x021 & FSR\_RAW\_8\_9 & uint16\_t{[}2{]} & RO & 0 & FSR sensors 8-9 \\
0x022 & FSR\_RAW\_10\_11 & uint16\_t{[}2{]} & RO & 0 & FSR sensors
10-11 \\
0x023 & FSR\_RAW\_12\_13 & uint16\_t{[}2{]} & RO & 0 & FSR sensors
12-13 \\
0x024 & FSR\_RAW\_14\_15 & uint16\_t{[}2{]} & RO & 0 & FSR sensors
14-15 \\
\end{longtable}
}

\textbf{FSR Layout:} 5 fingers × 6 FSRs/finger = 30 FSR sensors total

\textbf{Data Format:} 12-bit ADC values (0-4095). Higher values indicate
greater force.

*Availability:** Populated when reply variant 3 (TX3) is selected via
command header.

\paragraph{5.3.11 Status and Reply Information
(0x025-0x026)}\label{status-and-reply-information-0x025-0x026}

{\def\LTcaptype{none} % do not increment counter
\begin{longtable}[]{@{}
  >{\raggedright\arraybackslash}p{(\linewidth - 10\tabcolsep) * \real{0.1905}}
  >{\raggedright\arraybackslash}p{(\linewidth - 10\tabcolsep) * \real{0.2381}}
  >{\raggedright\arraybackslash}p{(\linewidth - 10\tabcolsep) * \real{0.0952}}
  >{\raggedright\arraybackslash}p{(\linewidth - 10\tabcolsep) * \real{0.1270}}
  >{\raggedright\arraybackslash}p{(\linewidth - 10\tabcolsep) * \real{0.1429}}
  >{\raggedright\arraybackslash}p{(\linewidth - 10\tabcolsep) * \real{0.2063}}@{}}
\toprule\noalign{}
\begin{minipage}[b]{\linewidth}\raggedright
Word Index
\end{minipage} & \begin{minipage}[b]{\linewidth}\raggedright
Register Name
\end{minipage} & \begin{minipage}[b]{\linewidth}\raggedright
Type
\end{minipage} & \begin{minipage}[b]{\linewidth}\raggedright
Access
\end{minipage} & \begin{minipage}[b]{\linewidth}\raggedright
Default
\end{minipage} & \begin{minipage}[b]{\linewidth}\raggedright
Description
\end{minipage} \\
\midrule\noalign{}
\endhead
\bottomrule\noalign{}
\endlastfoot
0x025 & HOT\_COLD\_BITMASK & uint32\_t & RO & 0 & Motor temperature
status bitmask. Bit N set = motor N warning \\
0x026 & FORMAT\_HEADER & uint32\_t & RO & 0 & Last received reply format
header. Indicates reply data structure \\
\end{longtable}
}

\textbf{HOT\_COLD\_BITMASK Bit Definitions:} - Bit 0: Motor 0
temperature warning - Bit 1: Motor 1 temperature warning - \ldots{} -
Bit 5: Motor 5 temperature warning - Bits {[}31:6{]}: Reserved

\textbf{FORMAT\_HEADER:} Echoes the command header from the most recent
valid reply. Use to verify expected reply variant was received.

\begin{center}\rule{0.5\linewidth}{0.5pt}\end{center}

\subsubsection{5.4 Block 4: UART Direct Buffers
(0x02A-0x051)}\label{block-4-uart-direct-buffers-0x02a-0x051}

This block provides low-level direct access to UART transmission and
reception buffers for pass-through mode. The controller is responsible
for HDLC byte stuffing on TX data. The firmware automatically performs
HDLC unstuffing on RX data.

\paragraph{5.4.1 Timeout Configuration
(0x02A)}\label{timeout-configuration-0x02a}

{\def\LTcaptype{none} % do not increment counter
\begin{longtable}[]{@{}
  >{\raggedright\arraybackslash}p{(\linewidth - 10\tabcolsep) * \real{0.1905}}
  >{\raggedright\arraybackslash}p{(\linewidth - 10\tabcolsep) * \real{0.2381}}
  >{\raggedright\arraybackslash}p{(\linewidth - 10\tabcolsep) * \real{0.0952}}
  >{\raggedright\arraybackslash}p{(\linewidth - 10\tabcolsep) * \real{0.1270}}
  >{\raggedright\arraybackslash}p{(\linewidth - 10\tabcolsep) * \real{0.1429}}
  >{\raggedright\arraybackslash}p{(\linewidth - 10\tabcolsep) * \real{0.2063}}@{}}
\toprule\noalign{}
\begin{minipage}[b]{\linewidth}\raggedright
Word Index
\end{minipage} & \begin{minipage}[b]{\linewidth}\raggedright
Register Name
\end{minipage} & \begin{minipage}[b]{\linewidth}\raggedright
Type
\end{minipage} & \begin{minipage}[b]{\linewidth}\raggedright
Access
\end{minipage} & \begin{minipage}[b]{\linewidth}\raggedright
Default
\end{minipage} & \begin{minipage}[b]{\linewidth}\raggedright
Description
\end{minipage} \\
\midrule\noalign{}
\endhead
\bottomrule\noalign{}
\endlastfoot
0x02A & ABH\_READ\_TIMEOUT & uint32\_t & R/W & 5 & Reply timeout in
milliseconds. Used in automatic mode \\
\end{longtable}
}

\textbf{Usage:} Defines how long the firmware waits for a UART reply
before marking the transaction as timed out.

\paragraph{5.4.2 UART RX Buffer
(0x02B-0x03E)}\label{uart-rx-buffer-0x02b-0x03e}

{\def\LTcaptype{none} % do not increment counter
\begin{longtable}[]{@{}
  >{\raggedright\arraybackslash}p{(\linewidth - 12\tabcolsep) * \real{0.2182}}
  >{\raggedright\arraybackslash}p{(\linewidth - 12\tabcolsep) * \real{0.2727}}
  >{\raggedright\arraybackslash}p{(\linewidth - 12\tabcolsep) * \real{0.1091}}
  >{\raggedright\arraybackslash}p{(\linewidth - 12\tabcolsep) * \real{0.1455}}
  >{\raggedright\arraybackslash}p{(\linewidth - 12\tabcolsep) * \real{0.0727}}
  >{\raggedright\arraybackslash}p{(\linewidth - 12\tabcolsep) * \real{0.0364}}
  >{\raggedright\arraybackslash}p{(\linewidth - 12\tabcolsep) * \real{0.1455}}@{}}
\toprule\noalign{}
\begin{minipage}[b]{\linewidth}\raggedright
Word Index
\end{minipage} & \begin{minipage}[b]{\linewidth}\raggedright
Register Name
\end{minipage} & \begin{minipage}[b]{\linewidth}\raggedright
Type
\end{minipage} & \begin{minipage}[b]{\linewidth}\raggedright
Access
\end{minipage} & \begin{minipage}[b]{\linewidth}\raggedright
Default
\end{minipage} & \begin{minipage}[b]{\linewidth}\raggedright
Size
\end{minipage} & \begin{minipage}[b]{\linewidth}\raggedright
Description
\end{minipage} \\
\midrule\noalign{}
\endhead
\bottomrule\noalign{}
\endlastfoot
0x02B-0x03D & UART\_RX\_DECODED{[}0-75{]} & uint8\_t{[}76{]} & RO & 0 &
19 words (76 bytes) & HDLC-decoded receive buffer. Contains unstuffed
UART data \\
0x03E & NBYTES\_DECODED\_UART & uint32\_t & RO & 0 & 1 word & Number of
valid bytes in UART\_RX\_DECODED buffer \\
\end{longtable}
}

\textbf{Buffer Layout:} 76-byte buffer. Firmware performs HDLC
unstuffing automatically, so this buffer always contains an unstuffed
payload when NBYTES\_DECODED\_UART is nonzero.

\paragraph{5.4.3 UART TX Buffer
(0x03F-0x052)}\label{uart-tx-buffer-0x03f-0x052}

{\def\LTcaptype{none} % do not increment counter
\begin{longtable}[]{@{}
  >{\raggedright\arraybackslash}p{(\linewidth - 12\tabcolsep) * \real{0.2000}}
  >{\raggedright\arraybackslash}p{(\linewidth - 12\tabcolsep) * \real{0.2500}}
  >{\raggedright\arraybackslash}p{(\linewidth - 12\tabcolsep) * \real{0.1000}}
  >{\raggedright\arraybackslash}p{(\linewidth - 12\tabcolsep) * \real{0.1333}}
  >{\raggedright\arraybackslash}p{(\linewidth - 12\tabcolsep) * \real{0.0500}}
  >{\raggedright\arraybackslash}p{(\linewidth - 12\tabcolsep) * \real{0.0500}}
  >{\raggedright\arraybackslash}p{(\linewidth - 12\tabcolsep) * \real{0.2167}}@{}}
\toprule\noalign{}
\begin{minipage}[b]{\linewidth}\raggedright
Word Index
\end{minipage} & \begin{minipage}[b]{\linewidth}\raggedright
Register Name
\end{minipage} & \begin{minipage}[b]{\linewidth}\raggedright
Type
\end{minipage} & \begin{minipage}[b]{\linewidth}\raggedright
Access
\end{minipage} & \begin{minipage}[b]{\linewidth}\raggedright
Default
\end{minipage} & \begin{minipage}[b]{\linewidth}\raggedright
Size
\end{minipage} & \begin{minipage}[b]{\linewidth}\raggedright
Description
\end{minipage} \\
\midrule\noalign{}
\endhead
\bottomrule\noalign{}
\endlastfoot
0x03F-0x051 & UART\_TX\_MEM{[}0-75{]} & uint8\_t{[}76{]} & R/W & 0 & 19
words & Raw transmit buffer. Controller must perform HDLC stuffing \\
0x052 & NBYTES\_WRITE\_UART & uint32\_t & WO & 0 & 1 word & TX trigger
register. Write byte count to initiate transmission \\
\end{longtable}
}

\textbf{Buffer Layout:} 76-byte buffer packed into 19 32-bit words.

\textbf{Usage Sequence (Pass-Through Mode):} 1. Construct UART frame in
local buffer 2. Perform HDLC byte stuffing (escape 0x7E and 0x7D) 3.
Write stuffed bytes to UART\_TX\_MEM (0x03F-0x051) 4. Write byte count
to NBYTES\_WRITE\_UART (0x052) 5. Transmission begins immediately

\begin{center}\rule{0.5\linewidth}{0.5pt}\end{center}

\subsubsection{5.6 Block 5: System Information
(0x053-0x057)}\label{block-5-system-information-0x053-0x057}

This block provides read-only system information and control flags for
firmware management.

\paragraph{5.6.1 Firmware Version
(0x053-0x056)}\label{firmware-version-0x053-0x056}

{\def\LTcaptype{none} % do not increment counter
\begin{longtable}[]{@{}
  >{\raggedright\arraybackslash}p{(\linewidth - 10\tabcolsep) * \real{0.2000}}
  >{\raggedright\arraybackslash}p{(\linewidth - 10\tabcolsep) * \real{0.2500}}
  >{\raggedright\arraybackslash}p{(\linewidth - 10\tabcolsep) * \real{0.1000}}
  >{\raggedright\arraybackslash}p{(\linewidth - 10\tabcolsep) * \real{0.1333}}
  >{\raggedright\arraybackslash}p{(\linewidth - 10\tabcolsep) * \real{0.1000}}
  >{\raggedright\arraybackslash}p{(\linewidth - 10\tabcolsep) * \real{0.2167}}@{}}
\toprule\noalign{}
\begin{minipage}[b]{\linewidth}\raggedright
Word Index
\end{minipage} & \begin{minipage}[b]{\linewidth}\raggedright
Register Name
\end{minipage} & \begin{minipage}[b]{\linewidth}\raggedright
Type
\end{minipage} & \begin{minipage}[b]{\linewidth}\raggedright
Access
\end{minipage} & \begin{minipage}[b]{\linewidth}\raggedright
Size
\end{minipage} & \begin{minipage}[b]{\linewidth}\raggedright
Description
\end{minipage} \\
\midrule\noalign{}
\endhead
\bottomrule\noalign{}
\endlastfoot
0x053-0x056 & GIT\_HASH\_BUFFER{[}0-15{]} & uint8\_t{[}16{]} & RO & 4
words & Git commit hash identifying firmware version (ASCII hex,
null-terminated, variable length) \\
\end{longtable}
}

\textbf{Format:} Null terminated ASCII string representing the --short
git commit SHA for the
\href{https://github.com/psyonicinc/fdcan-abh-converter-firmware}{interface
firmware repository}.

\textbf{Example:} ``38c0d64'' represents git commit hash starting with
38c0d64\ldots{}

\textbf{Usage:} Read to verify firmware version during system
initialization or diagnostics.

\paragraph{5.6.2 Control Flags (0x057)}\label{control-flags-0x057}

{\def\LTcaptype{none} % do not increment counter
\begin{longtable}[]{@{}
  >{\raggedright\arraybackslash}p{(\linewidth - 10\tabcolsep) * \real{0.1905}}
  >{\raggedright\arraybackslash}p{(\linewidth - 10\tabcolsep) * \real{0.2381}}
  >{\raggedright\arraybackslash}p{(\linewidth - 10\tabcolsep) * \real{0.0952}}
  >{\raggedright\arraybackslash}p{(\linewidth - 10\tabcolsep) * \real{0.1270}}
  >{\raggedright\arraybackslash}p{(\linewidth - 10\tabcolsep) * \real{0.1429}}
  >{\raggedright\arraybackslash}p{(\linewidth - 10\tabcolsep) * \real{0.2063}}@{}}
\toprule\noalign{}
\begin{minipage}[b]{\linewidth}\raggedright
Word Index
\end{minipage} & \begin{minipage}[b]{\linewidth}\raggedright
Register Name
\end{minipage} & \begin{minipage}[b]{\linewidth}\raggedright
Type
\end{minipage} & \begin{minipage}[b]{\linewidth}\raggedright
Access
\end{minipage} & \begin{minipage}[b]{\linewidth}\raggedright
Default
\end{minipage} & \begin{minipage}[b]{\linewidth}\raggedright
Description
\end{minipage} \\
\midrule\noalign{}
\endhead
\bottomrule\noalign{}
\endlastfoot
0x057 & UPDATE\_NONVOLATILE\_STORAGE & uint8\_t & R/W & 0 & Flash update
trigger. Write 1 to persist Block 1 registers. Auto-clears \\
\end{longtable}
}

\textbf{Usage Sequence:} 1. Modify one or more registers in Block 1
(0x000-0x005) 2. Write 1 to UPDATE\_NONVOLATILE\_STORAGE (0x057) 3.
Firmware commits changes to flash (Page 63) 4. Flag automatically clears
to 0 upon completion 5. Changes persist across power cycles

\textbf{Flash Write Time:} Approximately 100-200 ms (firmware blocks
during write)

\textbf{Caution:} Excessive flash writes reduce flash lifespan. Limit to
configuration changes only (not runtime updates).

\begin{center}\rule{0.5\linewidth}{0.5pt}\end{center}

\subsubsection{5.7 Register Access Notes}\label{register-access-notes}

\paragraph{5.7.1 Endianness}\label{endianness}

All multi-byte values use \textbf{little-endian} byte ordering (LSB at
lowest address).

\textbf{Example:} 32-bit value 0x12345678 stored as:

\begin{verbatim}
Byte 0: 0x78
Byte 1: 0x56
Byte 2: 0x34
Byte 3: 0x12
\end{verbatim}

\paragraph{5.7.2 Packed Arrays}\label{packed-arrays}

Registers containing arrays of int16\_t or uint16\_t pack two values per
32-bit word. Access pattern: - \textbf{Bits {[}15:0{]}:} Array element N
- \textbf{Bits {[}31:16{]}:} Array element N+1

\textbf{Example:} Q\_DESIRED\_0\_1 (0x008):

\begin{verbatim}
Bits [15:0]:  Motor 0 position setpoint (int16_t)
Bits [31:16]: Motor 1 position setpoint (int16_t)
\end{verbatim}

\paragraph{5.7.3 Atomic Access}\label{atomic-access}

The DARTT protocol ensures atomic 32-bit word access. Sub-word fields
(uint8\_t, uint16\_t) should be read/written as complete 32-bit words
with appropriate masking and shifting.

\paragraph{5.7.4 Write Triggers}\label{write-triggers}

Writing to \textbf{NBYTES\_WRITE\_UART} (0x052) triggers immediate UART
transmission. Ensure UART\_TX\_MEM is fully populated before triggering.

\paragraph{5.7.5 Read-Only Enforcement}\label{read-only-enforcement}

Read-only register designations are \textbf{informational only}. The
firmware does not enforce write protection. Writing to RO registers may
have undefined behavior.

\begin{center}\rule{0.5\linewidth}{0.5pt}\end{center}

\subsection{6. Operational Modes}\label{operational-modes}

The FDCAN-ABH Converter supports two distinct operational modes to
accommodate different levels of control abstraction.

\subsubsection{6.1 Mode 1: Automatic Ability Hand
Control}\label{mode-1-automatic-ability-hand-control}

In this mode, the converter handles all low-level protocol details,
including frame construction, checksum calculation, and HDLC byte
stuffing.

\paragraph{6.1.1 Characteristics}\label{characteristics}

\begin{itemize}
\tightlist
\item
  \textbf{User Interaction:} Read/write registers in Block 2 (Ability
  Hand Control API)
\item
  \textbf{Frame Generation:} Automatic
\item
  \textbf{HDLC Processing:} Automatic (both TX and RX)
\item
  \textbf{Checksum:} Calculated automatically
\item
  \textbf{Reply Parsing:} Automatic
\item
  \textbf{Use Case:} Standard motor control applications
\end{itemize}

\paragraph{6.1.2 Advantages}\label{advantages}

\begin{itemize}
\tightlist
\item
  \textbf{Minimizes latency}. Most suitable for high-frequency control
  loops
\item
  Simple controller implementation
\item
  Guaranteed protocol correctness
\item
  Automatic error handling
\end{itemize}

\paragraph{6.1.3 Limitations}\label{limitations}

\begin{itemize}
\tightlist
\item
  Limited to standard Ability Hand command set
\item
  Cannot implement custom protocols
\item
  Fixed frame structure
\end{itemize}

\subsubsection{6.2 Mode 2: Direct UART
Pass-Through}\label{mode-2-direct-uart-pass-through}

In this mode, the controller has direct access to UART buffers and is
responsible for frame construction and HDLC stuffing.

\paragraph{6.2.1 Characteristics}\label{characteristics-1}

\begin{itemize}
\tightlist
\item
  \textbf{User Interaction:} Read/write registers in Block 4 (UART
  Direct Buffers)
\item
  \textbf{Frame Generation:} Manual (controller responsibility)
\item
  \textbf{HDLC TX Stuffing:} Manual (controller responsibility)
\item
  \textbf{HDLC RX Unstuffing:} Automatic
\item
  \textbf{Checksum:} Manual (controller responsibility)
\item
  \textbf{Reply Parsing:} Manual (controller responsibility)
\item
  \textbf{Use Case:} Custom protocols, debugging, advanced features
\end{itemize}

\paragraph{6.2.2 Advantages}\label{advantages-1}

\begin{itemize}
\tightlist
\item
  Full protocol flexibility
\item
  Access to custom Ability Hand features
\item
  Debugging and diagnostic capabilities
\item
  Can implement non-standard command sequences
\end{itemize}

\paragraph{6.2.3 Limitations}\label{limitations-1}

\begin{itemize}
\tightlist
\item
  Controller must implement HDLC stuffing
\item
  Controller must implement checksum calculation
\item
  Higher controller complexity
\item
  Potential for protocol errors
\end{itemize}

\subsubsection{6.3 Mode Selection}\label{mode-selection}

The operational mode is controlled by the \textbf{ABH\_COMMAND\_HEADER}
register (0x007):

\begin{itemize}
\tightlist
\item
  \textbf{Automatic Mode:} Set ABH\_COMMAND\_HEADER to a valid command
\item
  \textbf{Pass-Through Mode:} Set ABH\_COMMAND\_HEADER to 0x00 (or any
  unmapped value)
\end{itemize}

\textbf{CRITICAL:} Before using pass-through mode (direct UART buffer
access), the ABH\_COMMAND\_HEADER register \textbf{must} be set to 0x00
or another unmapped value. Failure to do so will cause internal
collisions between automatic frame generation and manual buffer writes,
resulting in undefined behavior.

\textbf{Mode Transition Procedure:} 1. To enter pass-through mode: Write
0x00 to ABH\_COMMAND\_HEADER (0x007) 2. To enter automatic mode: Write
desired command header to ABH\_COMMAND\_HEADER (0x007)

\begin{center}\rule{0.5\linewidth}{0.5pt}\end{center}

\subsection{7. Non-Volatile Configuration
Management}\label{non-volatile-configuration-management}

\subsubsection{7.1 Configuration Persistence
Mechanism}\label{configuration-persistence-mechanism}

\paragraph{7.1.1 Write Sequence (Internal
Firmware)}\label{write-sequence-internal-firmware}

\begin{enumerate}
\def\labelenumi{\arabic{enumi}.}
\tightlist
\item
  Controller modifies registers in Block 1 (0x000-0x005)
\item
  Controller writes 1 to UPDATE\_NONVOLATILE\_STORAGE (0x057)
\item
  Firmware unlocks flash controller
\item
  Firmware erases Page 63
\item
  Firmware writes new configuration to Page 63
\item
  Firmware locks flash controller
\item
  Firmware clears UPDATE\_NONVOLATILE\_STORAGE (0x057)
\end{enumerate}

\textbf{Duration:} 100-200 ms (blocking operation)

\paragraph{7.1.2 Read Sequence
(Boot-Time)}\label{read-sequence-boot-time}

\begin{enumerate}
\def\labelenumi{\arabic{enumi}.}
\tightlist
\item
  Firmware reads Page 63 into RAM (fds\_params\_t structure)
\item
  Firmware validates configuration (basic sanity checks)
\item
  If valid: Apply configuration to peripherals (CAN, UART)
\item
  If invalid: Use defaults (MODULE\_NUMBER=0x50,
  UART\_BAUD\_RATE=460800, etc.)
\end{enumerate}

\textbf{Duration:} \textless1 ms (non-blocking during normal operation)

\subsubsection{7.2 Flash Write
Limitations}\label{flash-write-limitations}

\paragraph{7.2.1 Endurance}\label{endurance}

\textbf{STM32G4 Flash Endurance:} 10,000 erase/write cycles (minimum,
per datasheet)

\textbf{Practical Limits:} - Configuration changes: Unlimited
(infrequent) - Do NOT write on every control cycle - Do NOT write more
than 10,000 times over device lifetime

\textbf{Recommended Usage:} - Write only during initial configuration -
Write only when user explicitly changes settings - Avoid automated
runtime writes

\paragraph{7.2.2 Power Loss During Write}\label{power-loss-during-write}

\textbf{Risk:} If power is lost during flash write (Step 4-5 of Section
7.1.1), flash page may be corrupted.

\textbf{Mitigation:} 1. Firmware validates configuration on boot
(Section 7.1.2 Step 2) 2. If validation fails, firmware uses safe
defaults 3. Controller should avoid triggering writes during critical
operations

\textbf{Best Practice:} - Perform configuration writes during safe
periods (e.g., system startup) - Avoid writes during active prosthetic
use - Implement timeout monitoring (poll 0x057 for completion)

\subsubsection{7.3 Default Configuration}\label{default-configuration}

The firmware uses these defaults:

{\def\LTcaptype{none} % do not increment counter
\begin{longtable}[]{@{}lll@{}}
\toprule\noalign{}
Parameter & Default Value & Register \\
\midrule\noalign{}
\endhead
\bottomrule\noalign{}
\endlastfoot
MODULE\_NUMBER & 0x50 & 0x000 \\
UART\_BAUD\_RATE & 460800 & 0x001 \\
FDCAN\_NBRP & 2 & 0x002 \\
FDCAN\_NTSEG1 & 135 & 0x003 \\
FDCAN\_NTSEG2 & 34 & 0x004 \\
UNUSED\_ZEROPAD & 0 & 0x005 \\
\end{longtable}
}

\textbf{CAN Bit Rate:} 500 kbit/s (derived from NBRP, NTSEG1, NTSEG2)

\subsubsection{7.4 Configuration Validation
Warning}\label{configuration-validation-warning}

\textbf{WARNING:} The firmware does not perform validation of
configuration parameters on boot or during flash writes. Invalid
configuration parameters, especially FDCAN timing parameters (NBRP,
NTSEG1, NTSEG2), will cause the device to fail CAN bus initialization
and become unresponsive to CAN communication.

\textbf{Consequences of Invalid FDCAN Parameters:} - Device will become
inoperable over CAN - Recovery requires either: 1. Restoring valid
parameters via SWD programmer, OR 2. Full chip erase and firmware
re-flash via SWD programmer

\textbf{Recommendation:} Implement controller-side validation of the
full flash block any time a change is made before setting the update
flag.

\begin{center}\rule{0.5\linewidth}{0.5pt}\end{center}

\subsection{8. Physical and Electrical
Characteristics}\label{physical-and-electrical-characteristics}

\subsubsection{8.1 Microcontroller
Platform}\label{microcontroller-platform}

{\def\LTcaptype{none} % do not increment counter
\begin{longtable}[]{@{}ll@{}}
\toprule\noalign{}
Parameter & Value \\
\midrule\noalign{}
\endhead
\bottomrule\noalign{}
\endlastfoot
MCU & STM32G431xx \\
Core & ARM Cortex-M4F \\
Clock Frequency & 170 MHz (max) \\
Flash & 128 KB \\
RAM & 32 KB \\
\end{longtable}
}

\subsubsection{8.2 CAN Interface}\label{can-interface-1}

\textbf{Integrated Components:} - TLE9250XSJXUMA1 5V CAN transceiver
(on-board) - 120$\Omega$ termination resistor (on-board)

\textbf{Bit Rate:} 125 kbit/s to 1000 kbit/s (default: 500 kbit/s,
classic CAN compliant)

\textbf{External Connection:} - CANH and CANL differential pair to CAN
bus - Device acts as bus termination node (place at bus end or disable
external terminations)

\subsubsection{8.3 Power Supply}\label{power-supply}

{\def\LTcaptype{none} % do not increment counter
\begin{longtable}[]{@{}llllll@{}}
\toprule\noalign{}
Parameter & Min & Typ & Max & Unit & Notes \\
\midrule\noalign{}
\endhead
\bottomrule\noalign{}
\endlastfoot
Supply Voltage (VDD) & 2.0 & 8.4 & 12 & V & MCU supply \\
\end{longtable}
}

\begin{center}\rule{0.5\linewidth}{0.5pt}\end{center}

\subsection{9. Appendices}\label{appendices}

\begin{center}\rule{0.5\linewidth}{0.5pt}\end{center}

\subsubsection{Appendix A: Revision
History}\label{appendix-a-revision-history}

{\def\LTcaptype{none} % do not increment counter
\begin{longtable}[]{@{}llll@{}}
\toprule\noalign{}
Revision & Date & Author & Description \\
\midrule\noalign{}
\endhead
\bottomrule\noalign{}
\endlastfoot
1.0 & 2025-11-17 & Jesse Cornman & Initial release \\
\end{longtable}
}

\begin{center}\rule{0.5\linewidth}{0.5pt}\end{center}

\subsubsection{Appendix B: Glossary}\label{appendix-b-glossary}

{\def\LTcaptype{none} % do not increment counter
\begin{longtable}[]{@{}
  >{\raggedright\arraybackslash}p{(\linewidth - 2\tabcolsep) * \real{0.3333}}
  >{\raggedright\arraybackslash}p{(\linewidth - 2\tabcolsep) * \real{0.6667}}@{}}
\toprule\noalign{}
\begin{minipage}[b]{\linewidth}\raggedright
Term
\end{minipage} & \begin{minipage}[b]{\linewidth}\raggedright
Definition
\end{minipage} \\
\midrule\noalign{}
\endhead
\bottomrule\noalign{}
\endlastfoot
ABH & Ability Hand - Psyonic's prosthetic hand device \\
CAN & Controller Area Network - Serial communication protocol \\
DARTT & Dual Address Real-Time Transport - Block memory protocol over
CAN \\
FDCAN & Flexible Data-rate CAN - Enhanced CAN protocol (ISO 11898-1) \\
FSR & Force-Sensitive Resistor - Pressure sensor \\
HDLC & High-Level Data Link Control - Byte stuffing protocol (RFC
1662) \\
ICD & Interface Control Document - Specification document \\
LSB & Least Significant Bit/Byte \\
MSB & Most Significant Bit/Byte \\
NV & Non-Volatile - Data retained across power cycles \\
PPP & Point-to-Point Protocol - HDLC framing variant \\
RO & Read-Only - Register cannot be written \\
R/W & Read/Write - Register can be read and written \\
UART & Universal Asynchronous Receiver/Transmitter - Serial interface \\
WO & Write-Only - Writing triggers an action \\
\end{longtable}
}

\begin{center}\rule{0.5\linewidth}{0.5pt}\end{center}

\subsection{Document Approval}\label{document-approval}

{\def\LTcaptype{none} % do not increment counter
\begin{longtable}[]{@{}llll@{}}
\toprule\noalign{}
Role & Name & Signature & Date \\
\midrule\noalign{}
\endhead
\bottomrule\noalign{}
\endlastfoot
Author & Jesse Cornman & & 2025-11-17 \\
Technical Reviewer & & & \\
Project Manager & & & \\
\end{longtable}
}

\begin{center}\rule{0.5\linewidth}{0.5pt}\end{center}

\textbf{END OF DOCUMENT}

\end{document}
